\section{Hva er Internett?}
%=============================================================================

Internett kan beskrives fra to perspektiver: \textbf{``nuts and bolts''} (komponentene) og \textbf{``service''} (tjenestene).

\subsection{``Nuts and Bolts'' -- Komponentperspektivet}

\begin{keybox}[Internett er et ``network of networks'']
Internett best\aa r av milliarder tilkoblede enheter som kommuniserer via \textbf{pakkesvitsjer} (rutere og svitsjer) over \textbf{kommunikasjonslinker} (fiber, kobber, radio, satellitt).
\end{keybox}

\begin{center}
\begin{tikzpicture}[
    host/.style={rectangle, draw, fill=blue!10, minimum width=1.5cm, minimum height=0.8cm, font=\footnotesize},
    router/.style={circle, draw, fill=orange!30, minimum size=0.9cm, font=\footnotesize},
    cloud/.style={ellipse, draw, fill=blue!5, minimum width=4cm, minimum height=2cm},
    link/.style={thick, -},
    >=Stealth
]
    % ISP Cloud
    \node[cloud] (isp) at (0,0) {};
    \node[font=\footnotesize\bfseries, text=ntnu] at (0,0.3) {Nettverkskjerne};
    \node[font=\footnotesize, text=ntnu] at (0,-0.2) {(ISP-er, rutere)};

    % Routers inside
    \node[router] (r1) at (-0.8,0) {R};
    \node[router] (r2) at (0.8,0) {R};
    \draw[link] (r1) -- (r2);

    % Home network
    \node[host] (pc1) at (-4.5, 1.5) {PC};
    \node[host] (phone) at (-4.5, 0) {Mobil};
    \node[router] (hr) at (-3, 0.75) {\tiny WiFi};
    \draw[link] (pc1) -- (hr);
    \draw[link] (phone) -- (hr);
    \draw[link, blue!60, thick] (hr) -- (r1) node[midway, above, font=\tiny] {aksessnett};

    % Enterprise
    \node[host] (srv) at (4.5, 1.5) {Server};
    \node[host] (ws) at (4.5, 0) {PC};
    \node[router] (er) at (3, 0.75) {\tiny Sw};
    \draw[link] (srv) -- (er);
    \draw[link] (ws) -- (er);
    \draw[link, blue!60, thick] (er) -- (r2) node[midway, above, font=\tiny] {aksessnett};

    % Labels
    \node[font=\footnotesize\itshape, text=gray] at (-4.5, -1) {Hjemmenettverk};
    \node[font=\footnotesize\itshape, text=gray] at (4.5, -1) {Bedriftsnettverk};

    % Annotations
    \draw[decorate, decoration={brace,amplitude=5pt}, thick, ntnu] (-5.2,-0.5) -- (-5.2,2) node[midway, left, font=\footnotesize, text=ntnu, xshift=-3pt] {Edge};
    \draw[decorate, decoration={brace,amplitude=5pt}, thick, orange!70!black] (-1.8,-1) -- (1.8,-1) node[midway, below, font=\footnotesize, text=orange!70!black, yshift=-3pt] {Core};
\end{tikzpicture}
\end{center}

\textbf{Viktige komponenter:}
\begin{itemize}[noitemsep]
    \item \textbf{Hosts / endesystemer} -- enheter som kj\o rer nettverksapplikasjoner (PC-er, mobiler, servere, IoT)
    \item \textbf{Pakkesvitsjer} -- videresender pakker: \emph{rutere} (nettverkslaget) og \emph{svitsjer} (linklaget)
    \item \textbf{Kommunikasjonslinker} -- fiber, kobber, radio, satellitt. Overf\o ringsrate = \emph{b\aa ndbredde}
    \item \textbf{Nettverk} -- samling av enheter, rutere og linker styrt av en organisasjon
    \item \textbf{ISP} (Internet Service Provider) -- gir tilgang til Internett
\end{itemize}

\begin{keybox}[Internettstandarder]
\begin{itemize}[noitemsep]
    \item \textbf{RFC} (Request for Comments) -- dokumenter som beskriver standardene
    \item \textbf{IETF} (Internet Engineering Task Force) -- utvikler standardene
\end{itemize}
\end{keybox}

\subsection{``Service'' -- Tjenesteperspektivet}

Internett er en \textbf{infrastruktur som leverer tjenester til applikasjoner}: web, video, e-post, spill, e-handel, sosiale medier, IoT, etc.

Internett tilbyr et \textbf{programmeringsgrensesnitt} (socket API) til distribuerte applikasjoner -- ``hooks'' som lar apper koble seg til og bruke transporttjenester.

%=============================================================================
\section{Hva er en protokoll?}
%=============================================================================

\begin{keybox}[Definisjon: Protokoll]
En \textbf{protokoll} definerer \textbf{formatet}, \textbf{rekkef\o lgen} av meldinger som sendes og mottas mellom nettverksenheter, og \textbf{handlingene} som utf\o res n\aa r meldinger sendes eller mottas.
\end{keybox}

\begin{center}
\begin{tikzpicture}[>=Stealth, font=\footnotesize]
    % Human protocol
    \node[font=\bfseries, text=ntnu] at (-3.5, 4.5) {Menneskelig protokoll};
    \node[inner sep=0pt] (p1) at (-4.5, 4) {};
    \node[inner sep=0pt] (p2) at (-2.5, 4) {};

    \draw[->, thick, blue] (-4.3, 3.5) -- (-2.7, 3.0) node[midway, above, sloped] {Hei!};
    \draw[->, thick, blue] (-2.7, 2.5) -- (-4.3, 2.0) node[midway, above, sloped] {Hei!};
    \draw[->, thick, blue] (-4.3, 1.5) -- (-2.7, 1.0) node[midway, above, sloped] {Hva er klokka?};
    \draw[->, thick, blue] (-2.7, 0.5) -- (-4.3, 0.0) node[midway, above, sloped] {14:00};

    % Dashed line
    \draw[dashed, gray] (-4.5, 3.8) -- (-4.5, -0.3);
    \draw[dashed, gray] (-2.5, 3.8) -- (-2.5, -0.3);

    % Network protocol
    \node[font=\bfseries, text=ntnu] at (3.5, 4.5) {Nettverksprotokoll (TCP/HTTP)};
    \node[inner sep=0pt] (c) at (2, 4) {};
    \node[inner sep=0pt] (s) at (5, 4) {};

    \node[font=\tiny] at (2, 4.1) {Klient};
    \node[font=\tiny] at (5, 4.1) {Server};

    \draw[->, thick, red!70!black] (2.2, 3.3) -- (4.8, 2.8) node[midway, above, sloped] {TCP conn. request};
    \draw[->, thick, red!70!black] (4.8, 2.3) -- (2.2, 1.8) node[midway, above, sloped] {TCP conn. response};
    \draw[->, thick, red!70!black] (2.2, 1.3) -- (4.8, 0.8) node[midway, above, sloped] {GET /index.html};
    \draw[->, thick, red!70!black] (4.8, 0.3) -- (2.2, -0.2) node[midway, above, sloped] {<file>};

    \draw[dashed, gray] (2, 3.8) -- (2, -0.5);
    \draw[dashed, gray] (5, 3.8) -- (5, -0.5);

    % Time arrows
    \draw[->, gray, thick] (-4.8, 3.8) -- (-4.8, -0.3) node[below, font=\tiny] {tid};
    \draw[->, gray, thick] (1.7, 3.8) -- (1.7, -0.5) node[below, font=\tiny] {tid};
\end{tikzpicture}
\end{center}

All kommunikasjonsaktivitet p\aa\ Internett er styrt av protokoller.

%=============================================================================
\section{Nettverkets ytterkant (Network Edge)}
%=============================================================================

\subsection{Hosts: Klienter og servere}
\begin{itemize}[noitemsep]
    \item \textbf{Hosts} (endesystemer) sitter i nettverkets ``edge''
    \item To typer: \textbf{klienter} (foresp\o r tjenester) og \textbf{servere} (tilbyr tjenester)
    \item Servere er ofte i \textbf{datasentre}
\end{itemize}

\subsection{Aksessnettverk (Access Networks)}

Aksessnettverket kobler endesystemer til den f\o rste ruteren (``edge router'').

\bigskip
\begin{center}
\begin{tikzpicture}[
    box/.style={rectangle, draw, fill=#1, minimum width=3cm, minimum height=1.2cm, text width=2.8cm, align=center, font=\footnotesize},
    >=Stealth
]
    \node[box=blue!10] (dsl) at (0,0) {\textbf{DSL}\\Telefonlinje\\Ned: 24--52 Mbps\\Opp: 3.5--16 Mbps\\\textbf{Dedikert}};
    \node[box=green!10] (cable) at (4.5,0) {\textbf{Kabel (HFC)}\\TV-kabel\\Ned: 40 Mbps--1.2 Gbps\\Opp: 30--100 Mbps\\\textbf{Delt}};
    \node[box=orange!10] (fiber) at (9,0) {\textbf{Fiber (FTTH)}\\Fiberoptisk\\Opp til 10 Gbps\\\textbf{Dedikert}};

    \node[box=purple!10] (wifi) at (0,-2.5) {\textbf{WiFi (WLAN)}\\802.11b/g/n\\11, 54, 450 Mbps\\Rekkevidde: $\sim$30m};
    \node[box=red!10] (cell) at (4.5,-2.5) {\textbf{Mobilnett}\\4G/5G\\10--100+ Mbps\\Rekkevidde: $\sim$10 km};
    \node[box=cyan!10] (ent) at (9,-2.5) {\textbf{Enterprise}\\Ethernet + WiFi\\100 Mbps -- 10 Gbps};
\end{tikzpicture}
\end{center}

\begin{warnbox}[Delt vs. dedikert aksess]
\begin{itemize}[noitemsep]
    \item \textbf{DSL}: Dedikert linje til sentralen (DSLAM) -- ingen deling med naboer
    \item \textbf{Kabel}: Delt medium -- alle i nabolaget deler b\aa ndbredden til ``cable headend''
\end{itemize}
\end{warnbox}

\subsection{Hjemmenettverk}

\begin{center}
\begin{tikzpicture}[
    device/.style={rectangle, draw, fill=blue!5, minimum width=1.3cm, minimum height=0.6cm, font=\tiny},
    router/.style={rectangle, draw, fill=orange!20, minimum width=1.5cm, minimum height=0.6cm, font=\tiny},
    >=Stealth
]
    % Devices
    \node[device] (mob) at (-2, 1) {Mobil};
    \node[device] (pc) at (0, 1) {PC};
    \node[device] (tv) at (2, 1) {Smart-TV};

    % Router
    \node[router] (rtr) at (0, 0) {WiFi-ruter/NAT};

    % Modem
    \node[router] (mdm) at (0, -1) {Kabel/DSL-modem};

    % Links
    \draw[->, dashed] (mob) -- (rtr) node[midway, left, font=\tiny] {WiFi};
    \draw[->] (pc) -- (rtr) node[midway, right, font=\tiny] {Ethernet};
    \draw[->, dashed] (tv) -- (rtr) node[midway, right, font=\tiny] {WiFi};
    \draw[->] (rtr) -- (mdm);
    \draw[->, thick, red!60] (mdm) -- (0, -2) node[below, font=\tiny] {Til ISP / headend / sentral};
\end{tikzpicture}
\end{center}

\subsection{Hosts sender pakker}

\begin{formulabox}[Overf\o ringsforsinkelse (Transmission Delay)]
En host deler applikasjonsmeldinger i \textbf{pakker} med lengde $L$ bits og sender dem med overf\o ringsrate $R$ bits/sek:
\[
d_{\text{trans}} = \frac{L}{R}
\]
\end{formulabox}

\subsection{Fysiske medier}

\begin{center}
\begin{tabular}{lll}
\toprule
\textbf{Type} & \textbf{Kategori} & \textbf{Kjennetegn} \\
\midrule
Twisted pair (TP) & Guided & To isolerte kobberledere, Cat6: 10 Gbps \\
Koaksialkabel & Guided & To konsentriske ledere, bredb\aa nd \\
Fiberoptisk & Guided & Lyspulser i glass, h\o y hastighet, lav feil \\
Tr\aa dl\o s radio & Unguided & EM-b\o lger, p\aa virkes av refleksjon/st\o y \\
\bottomrule
\end{tabular}
\end{center}

\begin{itemize}[noitemsep]
    \item \textbf{Guided media}: Signal f\o lger fysisk medium (kobber, fiber, koaks)
    \item \textbf{Unguided media}: Signal g\aa r fritt (radio, WiFi, satellitt)
    \item Radio: kan v\ae re \textbf{half-duplex} (en retning om gangen)
\end{itemize}

%=============================================================================
\section{Nettverkskjernen (Network Core)}
%=============================================================================

Nettverkskjernen er et nettverk av sammenkoblede rutere -- et ``network of networks''.

\subsection{Pakkesvitsjing (Packet Switching)}

\begin{keybox}[Store-and-forward]
Hele pakken m\aa\ ankomme en ruter f\o r den kan videresendes til neste link.\\
Forsinkelse for \'{e}n pakke over \'{e}n link: $d = L/R$
\end{keybox}

To n\o kkelfunksjoner:
\begin{enumerate}[noitemsep]
    \item \textbf{Forwarding} (videresending): Flytte pakke fra inngangslink til riktig utgangslink
    \item \textbf{Routing} (ruting): Bestemme ruten pakken skal ta fra kilde til destinasjon
\end{enumerate}

\begin{center}
\begin{tikzpicture}[>=Stealth, font=\footnotesize]
    \node[rectangle, draw, fill=blue!10, minimum width=1.5cm] (src) at (0,0) {Kilde};
    \node[circle, draw, fill=orange!20, minimum size=1cm] (r1) at (3,0) {R1};
    \node[circle, draw, fill=orange!20, minimum size=1cm] (r2) at (6,0) {R2};
    \node[rectangle, draw, fill=green!10, minimum width=1.5cm] (dst) at (9,0) {Dest.};

    \draw[->, thick] (src) -- (r1) node[midway, above] {$L/R$};
    \draw[->, thick] (r1) -- (r2) node[midway, above] {$L/R$};
    \draw[->, thick] (r2) -- (dst) node[midway, above] {$L/R$};

    % Buffer illustration
    \draw[fill=yellow!20] (2.5, -0.8) rectangle (3.5, -1.2);
    \node[font=\tiny] at (3, -1) {buffer};
    \draw[->] (3, -0.5) -- (3, -0.8);

    \node[text=red, font=\tiny] at (3, -1.6) {K\o\ her $\rightarrow$ forsinkelse};
    \node[text=red, font=\tiny] at (3, -2.0) {Buffer full $\rightarrow$ pakketap};
\end{tikzpicture}
\end{center}

\textbf{K\o dannelse og pakketap:} Hvis pakker ankommer raskere enn de kan sendes videre, legges de i \textbf{buffer} (k\o). Hvis bufferen er full, \textbf{droppes} pakkene (tap).

\subsection{Kretssvitsjing (Circuit Switching)}

\begin{center}
\begin{tikzpicture}[font=\footnotesize]
    % Packet switching
    \node[font=\bfseries, text=ntnu] at (-3, 3) {Pakkesvitsjing};
    \draw[thick] (-4.5, 2.5) rectangle (-1.5, 0.5);
    \foreach \y in {2.3, 1.9, 1.5, 1.1, 0.7} {
        \draw[fill=blue!\y 0!white] (-4.3, \y) rectangle (-1.7, \y-0.25);
    }
    \node[font=\tiny, text width=2.5cm, align=center] at (-3, 0) {Delt link:\\pakker k\o er og deler};

    % Circuit switching
    \node[font=\bfseries, text=ntnu] at (3, 3) {Kretssvitsjing};
    \draw[thick] (1.5, 2.5) rectangle (4.5, 0.5);
    \draw[fill=red!30] (1.7, 2.3) rectangle (4.3, 1.9);
    \node[font=\tiny] at (3, 2.1) {Bruker 1 -- reservert};
    \draw[fill=green!30] (1.7, 1.7) rectangle (4.3, 1.3);
    \node[font=\tiny] at (3, 1.5) {Bruker 2 -- reservert};
    \draw[fill=gray!20] (1.7, 1.1) rectangle (4.3, 0.7);
    \node[font=\tiny] at (3, 0.9) {Ledig -- bortkastet};
    \node[font=\tiny, text width=2.5cm, align=center] at (3, 0) {Dedikert link:\\reservert per samtale};
\end{tikzpicture}
\end{center}

\begin{center}
\begin{tabular}{lcc}
\toprule
& \textbf{Pakkesvitsjing} & \textbf{Kretssvitsjing} \\
\midrule
Ressursdeling & Ja (statistisk multipleksing) & Nei (reservert) \\
Garanti & Ingen & Ja (fast b\aa ndbredde) \\
Effektivitet & H\o y (utnytter ledig kapasitet) & Lav (ledig = bortkastet) \\
K\o ing/tap & Ja & Nei \\
Best for & Bursty trafikk (web, e-post) & Konstant trafikk (telefon) \\
\bottomrule
\end{tabular}
\end{center}

\subsubsection{FDM og TDM}

\begin{center}
\begin{tikzpicture}[font=\footnotesize]
    % FDM
    \node[font=\bfseries, text=ntnu] at (-3, 2.5) {FDM};
    \draw[thick] (-5, 2) rectangle (-1, -0.5);
    \draw[fill=red!30] (-4.8, 1.8) rectangle (-1.2, 1.2);
    \node[font=\tiny] at (-3, 1.5) {Frekvensb\aa nd 1};
    \draw[fill=blue!30] (-4.8, 1.0) rectangle (-1.2, 0.4);
    \node[font=\tiny] at (-3, 0.7) {Frekvensb\aa nd 2};
    \draw[fill=green!30] (-4.8, 0.2) rectangle (-1.2, -0.3);
    \node[font=\tiny] at (-3, -0.05) {Frekvensb\aa nd 3};
    \node[font=\tiny, text width=3cm, align=center] at (-3, -1) {Hver bruker f\aa r\\eget frekvensb\aa nd};

    % TDM
    \node[font=\bfseries, text=ntnu] at (3, 2.5) {TDM};
    \draw[thick] (1, 2) rectangle (5, 0.5);
    \foreach \x/\col in {1.2/red, 2.0/blue, 2.8/green, 3.6/red} {
        \draw[fill=\col!30] (\x, 1.8) rectangle (\x+0.6, 0.7);
    }
    \draw[->] (1, 0.2) -- (5, 0.2) node[right, font=\tiny] {tid};
    \node[font=\tiny] at (1.5, 0.4) {t1};
    \node[font=\tiny] at (2.3, 0.4) {t2};
    \node[font=\tiny] at (3.1, 0.4) {t3};
    \node[font=\tiny] at (3.9, 0.4) {t4};
    \node[font=\tiny, text width=3cm, align=center] at (3, -0.5) {Hver bruker f\aa r\\egne tidsluker};
\end{tikzpicture}
\end{center}

\subsection{Internettstruktur}

\begin{center}
\begin{tikzpicture}[
    isp/.style={ellipse, draw, fill=#1, minimum width=2cm, minimum height=1cm, font=\footnotesize},
    >=Stealth
]
    % Tier 1
    \node[isp=red!20] (t1a) at (-1.5, 3) {Tier-1 ISP};
    \node[isp=red!20] (t1b) at (1.5, 3) {Tier-1 ISP};
    \draw[<->, thick] (t1a) -- (t1b);

    % IXP
    \node[rectangle, draw, fill=yellow!30, font=\footnotesize] (ixp) at (0, 2) {IXP};
    \draw[<->] (t1a) -- (ixp);
    \draw[<->] (t1b) -- (ixp);

    % Regional
    \node[isp=blue!15] (r1) at (-3, 1) {Regional ISP};
    \node[isp=blue!15] (r2) at (3, 1) {Regional ISP};
    \draw[->] (r1) -- (t1a);
    \draw[->] (r2) -- (t1b);

    % Access
    \node[isp=green!15] (a1) at (-4, -0.5) {\tiny Aksess-ISP};
    \node[isp=green!15] (a2) at (-2, -0.5) {\tiny Aksess-ISP};
    \node[isp=green!15] (a3) at (2, -0.5) {\tiny Aksess-ISP};
    \node[isp=green!15] (a4) at (4, -0.5) {\tiny Aksess-ISP};
    \draw[->] (a1) -- (r1);
    \draw[->] (a2) -- (r1);
    \draw[->] (a3) -- (r2);
    \draw[->] (a4) -- (r2);

    % Content provider
    \node[isp=purple!15] (goog) at (0, 0) {\tiny Google/Netflix};
    \draw[<->, dashed] (goog) -- (ixp);
    \draw[<->, dashed] (goog) -- (r1);

    % Labels
    \node[font=\tiny, text=gray] at (-4, -1.2) {brukere kobler seg til her};
\end{tikzpicture}
\end{center}

\begin{itemize}[noitemsep]
    \item \textbf{Aksess-ISP}: Gir brukere tilgang (Telenor, Altibox, etc.)
    \item \textbf{Regionale ISP}: Kobler aksess-ISP-er sammen
    \item \textbf{Tier-1 ISP}: Store, globale nettverk (f.eks.\ Telenor, AT\&T)
    \item \textbf{IXP} (Internet Exchange Point): Kobler ISP-er direkte
    \item \textbf{Content provider networks} (Google, Netflix): Bygger egne nettverk
\end{itemize}

%=============================================================================
\section{Ytelse: Forsinkelse, Tap og Throughput}
%=============================================================================

\subsection{Fire typer forsinkelse}

\begin{formulabox}[Total nodal delay]
\[
d_{\text{nodal}} = d_{\text{proc}} + d_{\text{queue}} + d_{\text{trans}} + d_{\text{prop}}
\]
\end{formulabox}

\begin{center}
\begin{tikzpicture}[>=Stealth, font=\footnotesize]
    % Router
    \node[circle, draw, fill=orange!20, minimum size=1.5cm] (r1) at (0,0) {Ruter};
    \node[circle, draw, fill=orange!20, minimum size=1.5cm] (r2) at (7,0) {Neste};

    % Packets in queue
    \foreach \x in {-1.5, -1.8, -2.1} {
        \draw[fill=blue!30] (\x, -0.15) rectangle (\x+0.25, 0.15);
    }

    % Packet on link
    \draw[fill=green!40] (3, -0.1) rectangle (3.3, 0.1);

    % Link
    \draw[thick] (0.75, 0) -- (6.25, 0);

    % Braces and labels
    \draw[decorate, decoration={brace,amplitude=4pt,mirror}] (-2.3, -0.5) -- (-1.3, -0.5)
        node[midway, below, yshift=-4pt, font=\tiny, text=red] {$d_{\text{queue}}$: k\o tid};

    \draw[decorate, decoration={brace,amplitude=4pt}] (-0.5, 0.9) -- (0.5, 0.9)
        node[midway, above, yshift=4pt, font=\tiny, text=red] {$d_{\text{proc}}$: prosessering};

    \draw[decorate, decoration={brace,amplitude=4pt,mirror}] (0.75, -0.5) -- (2.5, -0.5)
        node[midway, below, yshift=-4pt, font=\tiny, text=red] {$d_{\text{trans}} = L/R$};

    \draw[decorate, decoration={brace,amplitude=4pt}] (2.5, 0.5) -- (6.25, 0.5)
        node[midway, above, yshift=4pt, font=\tiny, text=red] {$d_{\text{prop}} = d/s$};
\end{tikzpicture}
\end{center}

\begin{center}
\begin{tabular}{llll}
\toprule
\textbf{Forsinkelse} & \textbf{Formel} & \textbf{Avhenger av} & \textbf{Typisk} \\
\midrule
$d_{\text{proc}}$ -- Prosessering & -- & Bitfeilsjekk, velge link & $< 1$ ms \\
$d_{\text{queue}}$ -- K\o & -- & Trafikkintensitet & varierer \\
$d_{\text{trans}}$ -- Overf\o ring & $L/R$ & Pakkest\o rrelse, linkrate & $\mu$s--ms \\
$d_{\text{prop}}$ -- Propagering & $d/s$ & Avstand, signalhastighet & $\mu$s--ms \\
\bottomrule
\end{tabular}
\end{center}

\begin{warnbox}[Viktig: Transmission delay $\neq$ Propagation delay!]
\begin{itemize}[noitemsep]
    \item $d_{\text{trans}} = L/R$ -- tid for \aa\ \textbf{dytte alle bits ut p\aa\ linken}
    \item $d_{\text{prop}} = d/s$ -- tid for signalet \aa\ \textbf{reise gjennom linken}
    \item $s \approx 2 \times 10^8$ m/s i fiber (lyshastighet / brytningsindeks 1.47)
\end{itemize}
\textbf{Eksempel:} Trondheim--Oslo (540 km), fiber:\\
$d_{\text{prop}} = 540\,000 / 200\,000\,000 = 2{,}7$ ms\\
Pakke 1500 bytes, link 1 Gbps:\\
$d_{\text{trans}} = 12\,000 / 1\,000\,000\,000 = 0{,}012$ ms
\end{warnbox}

\subsection{Trafikkintensitet og k\o forsinkelse}

\begin{formulabox}[Trafikkintensitet]
\[
\text{Trafikkintensitet} = \frac{L \cdot a}{R}
\]
hvor $L$ = pakkest\o rrelse (bits), $a$ = gjennomsnittlig ankomstrate (pakker/s), $R$ = linkrate (bps).
\end{formulabox}

\begin{center}
\begin{tikzpicture}[font=\footnotesize]
    \draw[->] (0,0) -- (6,0) node[right] {$La/R$};
    \draw[->] (0,0) -- (0,4.5) node[above] {Gjsn. k\o forsinkelse};

    % Curve
    \draw[thick, red, domain=0:0.95, samples=100] plot ({5*\x}, {0.3*\x/(1-\x)});

    % Asymptote
    \draw[dashed, gray] (5, 0) -- (5, 4.5);
    \node[below, font=\tiny] at (5, 0) {1};

    % Annotations
    \node[text=green!50!black, font=\tiny] at (1.5, 0.5) {$La/R \approx 0$: lite k\o};
    \node[text=orange!70!black, font=\tiny, text width=2cm, align=center] at (4, 2.5) {$La/R \to 1$:\\ stor k\o!};
    \node[text=red, font=\tiny, text width=2cm, align=center] at (6.5, 4) {$La/R > 1$:\\ uendelig!};
\end{tikzpicture}
\end{center}

\subsection{Pakketap}

N\aa r en pakke ankommer en ruter med \textbf{full buffer}, blir pakken \textbf{droppet} (tapt). Tapte pakker kan retransmitteres av forrige node, kilde-endesystemet, eller ikke i det hele tatt.

\subsection{Throughput}

\begin{keybox}[Throughput og flaskehals]
\textbf{Throughput}: Rate (bits/tidsenhet) som bits faktisk overf\o res fra sender til mottaker.
\begin{itemize}[noitemsep]
    \item \textbf{Instantaneous throughput}: Rate p\aa\ et gitt tidspunkt
    \item \textbf{Average throughput}: Rate over lengre tid
    \item \textbf{Bottleneck link}: Linken med lavest kapasitet bestemmer ende-til-ende throughput
\end{itemize}
\[
\text{Throughput} = \min(R_s, R_c)
\]
Med 10 samtidige forbindelser som deler en backbone $R$:
\[
\text{Per-connection throughput} = \min(R_c, R_s, R/10)
\]
\end{keybox}

\begin{center}
\begin{tikzpicture}[>=Stealth, font=\footnotesize]
    \node[rectangle, draw, fill=blue!10] (srv) at (0, 0) {Server};
    \node[circle, draw, fill=orange!20, minimum size=0.8cm] (r) at (4, 0) {R};
    \node[rectangle, draw, fill=green!10] (cli) at (8, 0) {Klient};

    % Pipes
    \draw[thick, fill=gray!10] (0.8, 0.3) -- (3.6, 0.3) -- (3.6, -0.3) -- (0.8, -0.3) -- cycle;
    \draw[thick, fill=gray!10] (4.4, 0.15) -- (7.2, 0.15) -- (7.2, -0.15) -- (4.4, -0.15) -- cycle;

    \node[font=\tiny, text=red] at (2.2, 0.5) {$R_s$ bits/sec};
    \node[font=\tiny, text=red] at (5.8, 0.35) {$R_c$ bits/sec};
    \node[font=\tiny] at (5.8, -0.5) {\textit{smalere = flaskehals}};
\end{tikzpicture}
\end{center}

\subsection{Traceroute}

\textbf{Traceroute} m\aa ler forsinkelse langs hele ruten fra kilde til destinasjon:
\begin{itemize}[noitemsep]
    \item Sender 3 pakker til hver ruter $i$ p\aa\ veien (med TTL = $i$)
    \item Ruter $i$ returnerer pakken, og senderen m\aa ler tiden
    \item Viser ogs\aa\ pakketap (``* * *'' = ingen respons)
\end{itemize}

%=============================================================================
\section{Sikkerhet}
%=============================================================================

Internett ble opprinnelig designet \textbf{uten} mye sikkerhet -- ``a group of mutually trusting users.''

\begin{center}
\begin{tabular}{lp{7cm}}
\toprule
\textbf{Trussel} & \textbf{Beskrivelse} \\
\midrule
\textbf{Malware} & Virus (krever brukeraksjon), ormer (selvreplikerende), spyware, botnets \\
\textbf{DoS/DDoS} & Overbelaster ressurser med falsk trafikk fra mange kompromitterte maskiner \\
\textbf{Packet sniffing} & Leser alle pakker p\aa\ delt medium (WiFi, Ethernet) med promiscuous NIC \\
\textbf{IP spoofing} & Sender pakker med falsk avsenderadresse \\
\bottomrule
\end{tabular}
\end{center}

\begin{keybox}[Forsvar]
\begin{itemize}[noitemsep]
    \item \textbf{Autentisering} -- beskytter mot spoofing
    \item \textbf{Kryptering} -- beskytter mot sniffing (konfidensialitet)
    \item \textbf{Digitale signaturer} -- sikrer integritet
    \item \textbf{Brannmurer} -- blokkerer uautorisert trafikk
    \item \textbf{VPN} -- begrenser tilgang til ressurser
\end{itemize}
\end{keybox}

%=============================================================================
\section{Protokollag og innkapsling}
%=============================================================================

\subsection{Internettets 5-lags protokollstakk}

\begin{center}
\begin{tikzpicture}[
    layer/.style={rectangle, draw, fill=#1, minimum width=7cm, minimum height=1cm, font=\footnotesize},
    >=Stealth
]
    \node[layer=red!15] (app) at (0, 4) {\textbf{5. Applikasjonslaget} -- HTTP, DNS, SMTP, FTP};
    \node[layer=orange!15] (trans) at (0, 3) {\textbf{4. Transportlaget} -- TCP, UDP};
    \node[layer=yellow!15] (net) at (0, 2) {\textbf{3. Nettverkslaget} -- IP, rutingprotokoller};
    \node[layer=green!15] (link) at (0, 1) {\textbf{2. Linklaget} -- Ethernet, WiFi};
    \node[layer=blue!15] (phys) at (0, 0) {\textbf{1. Fysisk lag} -- bits p\aa\ ledningen};

    % Descriptions on right
    \node[font=\tiny, text=gray, text width=4cm, right] at (3.8, 4) {St\o tter nettverksapper\\Data: \textbf{melding}};
    \node[font=\tiny, text=gray, text width=4cm, right] at (3.8, 3) {Prosess-til-prosess\\Data: \textbf{segment}};
    \node[font=\tiny, text=gray, text width=4cm, right] at (3.8, 2) {Host-til-host ruting\\Data: \textbf{datagram}};
    \node[font=\tiny, text=gray, text width=4cm, right] at (3.8, 1) {Nabo-til-nabo\\Data: \textbf{ramme (frame)}};
    \node[font=\tiny, text=gray, text width=4cm, right] at (3.8, 0) {Bits over medium};
\end{tikzpicture}
\end{center}

\begin{warnbox}[ISO/OSI-modellen har 7 lag]
OSI legger til \textbf{Presentasjon} og \textbf{Sesjon} mellom applikasjon og transport. Disse er i praksis inkludert i applikasjonslaget i Internett-modellen.
\end{warnbox}

\subsection{Innkapsling (Encapsulation)}

\begin{center}
\begin{tikzpicture}[font=\footnotesize, >=Stealth]
    % Application
    \draw[fill=red!15] (0, 4) rectangle (4, 4.5);
    \node at (2, 4.25) {Melding (M)};

    % Transport
    \draw[fill=orange!15] (-0.5, 3) rectangle (0, 3.5);
    \draw[fill=red!15] (0, 3) rectangle (4, 3.5);
    \node[font=\tiny] at (-0.25, 3.25) {$H_t$};
    \node at (2, 3.25) {Segment ($H_t$ + M)};

    % Network
    \draw[fill=yellow!15] (-1, 2) rectangle (-0.5, 2.5);
    \draw[fill=orange!15] (-0.5, 2) rectangle (0, 2.5);
    \draw[fill=red!15] (0, 2) rectangle (4, 2.5);
    \node[font=\tiny] at (-0.75, 2.25) {$H_n$};
    \node at (2, 2.25) {Datagram ($H_n$ + $H_t$ + M)};

    % Link
    \draw[fill=green!15] (-1.5, 1) rectangle (-1, 1.5);
    \draw[fill=yellow!15] (-1, 1) rectangle (-0.5, 1.5);
    \draw[fill=orange!15] (-0.5, 1) rectangle (0, 1.5);
    \draw[fill=red!15] (0, 1) rectangle (4, 1.5);
    \node[font=\tiny] at (-1.25, 1.25) {$H_l$};
    \node at (2, 1.25) {Ramme ($H_l$ + $H_n$ + $H_t$ + M)};

    % Arrows
    \draw[->, thick, gray] (4.5, 4.25) -- (4.5, 3.25);
    \draw[->, thick, gray] (4.5, 3.25) -- (4.5, 2.25);
    \draw[->, thick, gray] (4.5, 2.25) -- (4.5, 1.25);

    \node[font=\tiny, text=gray, rotate=-90] at (5, 2.75) {innkapsling};
\end{tikzpicture}
\end{center}

Hvert lag legger til sin \textbf{header} ($H$) rundt dataene fra laget over. Mottaker fjerner headers i omvendt rekkef\o lge.

\bigskip
\begin{keybox}[Hvorfor lagdeling?]
\begin{itemize}[noitemsep]
    \item Gir \textbf{klar struktur} -- lettere \aa\ designe og vedlikeholde
    \item \textbf{Modularisering} -- kan endre ett lag uten \aa\ p\aa virke andre
    \item Separerer komplekse systemer i h\aa ndterbare deler
\end{itemize}
\end{keybox}

%=============================================================================
\section{Pensumssp\o rsm\aa l -- Kapittel 1}
%=============================================================================

Her er de viktigste sp\o rsm\aa lene fra pensum, organisert etter tema.

\subsection{1.1 Nuts and bolts vs. Service view}

\begin{exambox}[Eksamensssp\o rsm\aa l]
\begin{longtable}{p{6.5cm}p{7cm}}
\toprule
\textbf{Sp\o rsm\aa l} & \textbf{Svar} \\
\midrule
\endhead
Hva er internett if\o lge ``nuts and bolts''? & Et globalt nettverk som kobler sammen milliarder av endesystemer. \\
\midrule
Hva er en ISP? & En internettleverand\o r som gir endesystemer tilgang til internett. \\
\midrule
Hvilke protokoller styrer internett? & TCP/IP. \\
\midrule
Hvem utvikler internettstandardene? & IETF (Internet Engineering Task Force). \\
\midrule
Hva kalles dokumentene som beskriver standardene? & RFC-er (Requests for Comments). \\
\midrule
Hva er internett if\o lge ``services view''? & En plattform som leverer tjenester til applikasjoner. \\
\midrule
Hva er en protokoll? & Et regelsett som definerer format, rekkef\o lge og handlinger for meldinger mellom kommunikasjonsenheter. \\
\bottomrule
\end{longtable}
\end{exambox}

\subsection{1.2 Network Edge}

\begin{exambox}[Eksamensssp\o rsm\aa l]
\begin{longtable}{p{6.5cm}p{7cm}}
\toprule
\textbf{Sp\o rsm\aa l} & \textbf{Svar} \\
\midrule
\endhead
Hvilke to typer verter finnes? & Klienter og servere. \\
\midrule
Hva st\aa r DSL for? & Digital Subscriber Line. \\
\midrule
Er DSL symmetrisk eller asymmetrisk? & Asymmetrisk (h\o yere nedlasting). \\
\midrule
Hva er HFC? & Hybrid Fiber Coax -- kabelnett med fiber + koaks. \\
\midrule
Hva er forskjellen p\aa\ delt og dedikert aksess? & DSL er dedikert, kabel er delt. \\
\midrule
Hva er formelen for packet transmission delay? & $D = L / R$. \\
\midrule
Hva er guided vs.\ unguided media? & Guided: signal i fysisk medium (kobber, fiber). Unguided: signal fritt (radio). \\
\midrule
Hva er fordelene med fiber? & H\o y hastighet, lav feilrate. \\
\bottomrule
\end{longtable}
\end{exambox}

\subsection{1.3 Network Core}

\begin{exambox}[Eksamensssp\o rsm\aa l]
\begin{longtable}{p{6.5cm}p{7cm}}
\toprule
\textbf{Sp\o rsm\aa l} & \textbf{Svar} \\
\midrule
\endhead
Hva er ``store-and-forward''? & Hele pakken m\aa\ ankomme f\o r den kan sendes videre. \\
\midrule
Hva skjer n\aa r bufferen er full? & Pakker droppes (tap). \\
\midrule
Hva er forskjellen p\aa\ pakkesvitsjing og kretssvitsjing? & Krets reserverer fast forbindelse; pakke sender data i sm\aa\ biter som deler linjen. \\
\midrule
Hva er FDM? & Frequency Division Multiplexing -- hver bruker f\aa r eget frekvensb\aa nd. \\
\midrule
Hva er TDM? & Time Division Multiplexing -- hver bruker f\aa r egne tidsluker. \\
\midrule
N\aa r er pakkesvitsjing best? & For bursty trafikk -- noen ganger aktiv, noen ganger ikke. \\
\midrule
Hva er tier-1 ISP? & Store, globale nettverk (f.eks.\ Telenor, AT\&T). \\
\midrule
Hva er IXP? & Internet Exchange Point -- kobler ISP-er direkte. \\
\bottomrule
\end{longtable}
\end{exambox}

\subsection{1.4 Performance}

\begin{exambox}[Eksamensssp\o rsm\aa l]
\begin{longtable}{p{6.5cm}p{7cm}}
\toprule
\textbf{Sp\o rsm\aa l} & \textbf{Svar} \\
\midrule
\endhead
Hva er de fire forsinkelsestypene? & Processing, queueing, transmission ($L/R$), propagation ($d/s$). \\
\midrule
Hva er trafikkintensitet? & $(L \times a) / R$ -- forholdet mellom ankomstrate og tjenesterate. \\
\midrule
Hva skjer n\aa r $La/R > 1$? & Mer arbeid ankommer enn kan h\aa ndteres; forsinkelse $\to \infty$. \\
\midrule
Hva er throughput? & Rate bits faktisk overf\o res fra sender til mottaker. \\
\midrule
Hva er bottleneck link? & Linken med lavest kapasitet -- bestemmer ende-til-ende throughput. \\
\bottomrule
\end{longtable}
\end{exambox}

\subsection{1.5 Lagdeling og 1.6 Sikkerhet}

\begin{exambox}[Eksamensssp\o rsm\aa l]
\begin{longtable}{p{6.5cm}p{7cm}}
\toprule
\textbf{Sp\o rsm\aa l} & \textbf{Svar} \\
\midrule
\endhead
Hvor mange lag i internett-protokollstakken? & Fem: applikasjon, transport, nettverk, link, fysisk. \\
\midrule
Hva er innkapsling? & Pakke data med headers fra hvert lag. \\
\midrule
Hva er packet sniffing? & Lese alle pakker p\aa\ delt nettverk. \\
\midrule
Hva er IP spoofing? & Sende pakker med falsk avsenderadresse. \\
\midrule
Hva er et DoS-angrep? & Gj\o re ressurser utilgjengelige ved \aa\ overbelaste med falsk trafikk. \\
\midrule
Hva beskytter mot sniffing? & Kryptering. \\
\midrule
Hva beskytter mot spoofing? & Autentisering. \\
\bottomrule
\end{longtable}
\end{exambox}

%=============================================================================
\section{Oppsummering -- viktige formler}
%=============================================================================

\begin{formulabox}[Formler du m\aa\ kunne]
\begin{align}
d_{\text{trans}} &= \frac{L}{R} & &\text{Transmission delay (pakkest\o rrelse / linkrate)} \\[6pt]
d_{\text{prop}} &= \frac{d}{s} & &\text{Propagation delay (avstand / signalhastighet)} \\[6pt]
d_{\text{nodal}} &= d_{\text{proc}} + d_{\text{queue}} + d_{\text{trans}} + d_{\text{prop}} & &\text{Total forsinkelse per node} \\[6pt]
\text{Trafikk\-intensitet} &= \frac{L \cdot a}{R} & &\text{Ankomstrate vs.\ tjenesterate} \\[6pt]
\text{Throughput} &= \min(R_s, R_c) & &\text{Bestemt av flaskehalsen}
\end{align}

\textbf{Verdier \aa\ huske:}
\begin{itemize}[noitemsep]
    \item Signalhastighet i fiber: $s \approx 2 \times 10^8$ m/s
    \item Lyshastighet: $c = 3 \times 10^8$ m/s (fiber har brytningsindeks $\approx 1{,}47$)
    \item 1 byte = 8 bits
\end{itemize}
\end{formulabox}
